\subsection{Matriz de requisitos funcionales y responsabilidad modular}

El Cuadro~\ref{tab:matriz_rf} relaciona los siete requisitos funcionales con los módulos responsables. Una marca indica responsabilidad de diseño, no grado de implementación ni resultado de verificación.

\begin{table}[H]
\caption{Trazabilidad entre requisitos funcionales y módulos}
\label{tab:matriz_rf}
\centering
\begin{tabularx}{\textwidth}{|X|c|c|c|c|c|c|}
\hline
\textbf{Requisito} & \textbf{M1} & \textbf{M2} & \textbf{M3} & \textbf{M4} & \textbf{M5} & \textbf{M6} \\
\hline
RF-01: Integración con fuentes de información hospitalaria & \checkmark & \checkmark & & & \checkmark & \checkmark \\
\hline
RF-02: Consulta del contexto clínico & & \checkmark & & & \checkmark & \checkmark \\
\hline
RF-03: Visualización de la bitácora asistencial & & \checkmark & & & \checkmark & \checkmark \\
\hline
RF-04: Registro de notas de contexto & & \checkmark & \checkmark & & \checkmark & \checkmark \\
\hline
RF-05: Transcripción de dictados & & & \checkmark & & & \checkmark \\
\hline
RF-06: Asistencia para la redacción & & & & \checkmark & & \checkmark \\
\hline
RF-07: Resúmenes asistidos & & \checkmark & & \checkmark & \checkmark & \checkmark \\
\hline
\end{tabularx}
\end{table}

M5 y M6 son transversales a las operaciones protegidas y al registro de procedencia y autoría. La matriz expresa la distribución prevista de responsabilidades; la evidencia de implementación y verificación se informa por separado.
